\documentclass[11pt,a4paper,twoside]{article}

% ── Overleaf-compatible preamble ──────────────────────────────────────────────
\usepackage[utf8]{inputenc}
\usepackage[T1]{fontenc}
\usepackage{amsmath,amssymb,amsthm}
\usepackage{mathtools}
\usepackage{xcolor}
\usepackage{hyperref}
\usepackage{enumitem}
\usepackage{booktabs}
\usepackage[margin=1in]{geometry}
\usepackage{fancyhdr}
\usepackage{lastpage}
\usepackage{tikz}
\usepackage{float}
\usepackage{array}
\usepackage{longtable}
\usepackage{multicol}

% ── Theorem environments ──────────────────────────────────────────────────────
\theoremstyle{definition}
\newtheorem{definition}{Definition}[section]
\newtheorem{rule}{Rule}[section]
\newtheorem{example}{Example}[section]
\newtheorem{proposition}{Proposition}[section]
\newtheorem{invariant}{Invariant}[section]

\theoremstyle{remark}
\newtheorem{remark}{Remark}[section]
\newtheorem{note}{Note}[section]

% ── Card notation ─────────────────────────────────────────────────────────────
\newcommand{\Card}[2]{\ensuremath{\langle #1, #2 \rangle}}
\newcommand{\Suit}[1]{\ensuremath{\text{#1}}}
\newcommand{\Rank}[1]{\ensuremath{\text{#1}}}
\newcommand{\Joker}[1]{\ensuremath{\mathcal{J}_{#1}}}
\newcommand{\Deck}{\ensuremath{\mathcal{D}}}
\newcommand{\Hand}{\ensuremath{\mathcal{H}}}
\newcommand{\Pile}{\ensuremath{\mathcal{P}}}
\newcommand{\State}{\ensuremath{\mathcal{S}}}
\newcommand{\Action}{\ensuremath{\mathcal{A}}}
\newcommand{\Top}{\ensuremath{\tau}}
\newcommand{\Dir}{\ensuremath{\delta}}
\newcommand{\Players}{\ensuremath{\Pi}}
\newcommand{\Winner}{\ensuremath{\mathcal{W}}}
\newcommand{\Cardless}{\ensuremath{\mathcal{C}}}
\newcommand{\Bomb}{\ensuremath{\mathcal{B}}}
\newcommand{\Quest}{\ensuremath{\mathcal{Q}}}
\newcommand{\OnVerge}{\ensuremath{\mathcal{V}}}
\newcommand{\Failed}{\ensuremath{\mathcal{F}}}
\newcommand{\CardReq}{\ensuremath{\mathcal{R}}}
\newcommand{\Special}{\ensuremath{\mathcal{S}}}
\newcommand{\Normal}{\ensuremath{\mathcal{N}}}
\newcommand{\func}[1]{\ensuremath{\mathsf{#1}}}

% ── Header / footer ───────────────────────────────────────────────────────────
\pagestyle{fancy}
\fancyhf{}
\fancyhead[L]{\small Kenyan Poker Rules}
\fancyhead[R]{\small v1.0}
\fancyfoot[C]{\thepage\ of \pageref{LastPage}}

% ── Title ─────────────────────────────────────────────────────────────────────
\title{\textbf{Kenyan Poker}\\[4pt]
\large A Formal Specification of Rules and Mechanics}
\author{}
\date{July 2026}

\begin{document}
\maketitle
\thispagestyle{empty}
\vspace{2cm}
\begin{center}
\emph{``Card!''}
\end{center}
\newpage

% ══════════════════════════════════════════════════════════════════════════════
\section{Introduction}
% ══════════════════════════════════════════════════════════════════════════════

Kenyan Poker is a shedding-type card game for 2 or more players, played with
one or more standard 54-card decks (52 standard cards plus 2 jokers). The
objective is to be the first player to empty one's hand by playing a
non-special card as the final card.

The game combines elements of Crazy Eights, UNO, and traditional accumulating
card games. Its defining features are a rich system of \emph{special cards}
that include cumulative bombs, question chains, jumps, kickbacks, and a
dual-power Ace of Spades, together with a ``declare card'' rule that punishes
players who fail to announce that they are on the verge of winning.

This document provides:
\begin{itemize}[nosep]
  \item A complete natural-language rulebook (Sections~\ref{sec:components}--\ref{sec:win}).
  \item A formal specification of the game state and transition system (Section~\ref{sec:formal}).
  \item Key invariants and properties (Section~\ref{sec:invariants}).
\end{itemize}

% ══════════════════════════════════════════════════════════════════════════════
\section{Components}
\label{sec:components}
% ══════════════════════════════════════════════════════════════════════════════

\subsection{The Deck}

One standard deck consists of 54 cards:
\begin{itemize}[nosep]
  \item 52 \emph{regular} cards: 4 suits $\times$ 13 ranks.
  \item 2 \emph{jokers}: one red ($\Joker{red}$), one black ($\Joker{black}$).
\end{itemize}

\begin{definition}[Card]
  A card $c$ is either a \textbf{regular card} $\Card{s}{r}$ where
  $s \in \{\heartsuit, \diamondsuit, \clubsuit, \spadesuit\}$ is a suit and
  $r \in \{A,2,3,4,5,6,7,8,9,10,J,Q,K\}$ is a rank, or a
  \textbf{joker} $\Joker{col}$ where $col \in \{\text{red}, \text{black}\}$.
\end{definition}

\begin{definition}[Card color]
  The color of a card, $\func{color}(c)$, is:
  \[
  \func{color}(c) = \begin{cases}
    \text{red}   & \text{if } c = \Card{s}{r} \text{ and } s \in \{\heartsuit, \diamondsuit\} \\
    \text{red}   & \text{if } c = \Joker{red} \\
    \text{black} & \text{if } c = \Card{s}{r} \text{ and } s \in \{\spadesuit, \clubsuit\} \\
    \text{black} & \text{if } c = \Joker{black}
  \end{cases}
  \]
\end{definition}

\subsection{Card Classification}

\begin{definition}[Special card]
  A card $c$ is \textbf{special}, denoted $c \in \Special$, iff:
  \begin{itemize}[nosep]
    \item $c$ is a joker, or
    \item $c = \Card{s}{r}$ and $r \in \{A, 2, 3, 8, J, Q, K\}$.
  \end{itemize}
  All other cards are \textbf{non-special} (or \emph{normal}), denoted $\Normal = \Deck \setminus \Special$.
\end{definition}

\begin{center}
\begin{tabular}{ccl}
\toprule
\textbf{Ranks} & \textbf{Classification} & \textbf{Special Effect} \\
\midrule
A    & Special   & Ace --- stops bombs; requests suit or number \\
2    & Special   & Bomb (2 cards) \\
3    & Special   & Bomb (3 cards) \\
4--7 & Normal    & --- \\
8    & Special   & Question \\
9,10 & Normal   & --- \\
J    & Special   & Jump --- skip next player \\
Q    & Special   & Question \\
K    & Special   & Kickback --- reverse direction \\
Joker& Special   & Bomb (5 cards); matches by color \\
\bottomrule
\end{tabular}
\end{center}

\subsection{Multiple Decks}

When the number of players $n$ exceeds a threshold $T$ (default $T = 6$), the
game uses $k = \lceil n / T \rceil$ standard decks, all shuffled together into
a single draw pile. The number of decks may also be set explicitly for a
faster or slower game.

% ══════════════════════════════════════════════════════════════════════════════
\section{Game Setup}
\label{sec:setup}
% ══════════════════════════════════════════════════════════════════════════════

\begin{enumerate}[label=\textbf{S\arabic*.}, leftmargin=*]

  \item \textbf{Shuffle.} Combine $k$ standard 54-card decks and shuffle
  thoroughly to form the \emph{draw pile} $\Pile_{\text{draw}}$.

  \item \textbf{Deal.} Deal $h$ cards to each player (default $h = 5$).
  Players may look at their own cards but keep them hidden from others.

  \item \textbf{Starter card.} Draw cards from the top of $\Pile_{\text{draw}}$
  one at a time, placing each into the \emph{discard pile}
  $\Pile_{\text{discard}}$, until a \textbf{non-special} card is drawn.
  This card becomes the \emph{top card} $\Top$, face-up, and the game begins.

  \item \textbf{Starting player.} Choose a starting player uniformly at random
  from the set of all players $\Players = \{0, 1, \dots, n-1\}$.

  \item \textbf{Direction.} The initial direction of play is clockwise
  ($\Dir = +1$).

\end{enumerate}

% ══════════════════════════════════════════════════════════════════════════════
\section{Core Mechanics}
\label{sec:core}
% ══════════════════════════════════════════════════════════════════════════════

\subsection{Legal Play}

\begin{definition}[Legal card]
  A card $c$ is \textbf{legal} to play on top of $\Top$ iff one of the
  following holds:
  \begin{enumerate}[nosep, label=(\roman*)]
    \item $c$ and $\Top$ are both regular and share a suit or rank:
          $c.\text{suit} = \Top.\text{suit} \;\lor\; c.\text{rank} = \Top.\text{rank}$.
    \item $c$ is a joker and $\func{color}(c) = \func{color}(\Top)$.
    \item $\Top$ is a joker and $\func{color}(c) = \func{color}(\Top)$.
  \end{enumerate}
\end{definition}

\begin{example}
  If $\Top = \Card{\spadesuit}{7}$, then $\Card{\spadesuit}{K}$ (same suit),
  $\Card{\heartsuit}{7}$ (same rank), and $\Joker{black}$ (matching color) are
  all legal. $\Card{\diamondsuit}{3}$ is \emph{not} legal.
\end{example}

\subsection{Turn Structure}

Play proceeds in turns around the table according to the current direction
$\Dir \in \{+1, -1\}$ (clockwise or counterclockwise). On a player's turn,
they must resolve any pending states in the following priority order:

\begin{enumerate}[label=\textbf{T\arabic*.}, leftmargin=*]
  \item \textbf{Cardless check.} If the player is \emph{cardless}
  (see~\S\ref{sec:cardless}), they must draw one card from $\Pile_{\text{draw}}$
  and their turn ends immediately.
  \item \textbf{Bomb response.} If a bomb is active targeting this player
  (see~\S\ref{sec:bomb}), they must respond.
  \item \textbf{Question response.} If a question chain is active and the
  player is the questioner (see~\S\ref{sec:question}), they must answer.
  \item \textbf{Ace request response.} If an Ace request is active
  (see~\S\ref{sec:ace}), they must honor it or draw.
  \item \textbf{Normal play.} The player must either play a legal card from
  their hand, or draw one card from $\Pile_{\text{draw}}$.
\end{enumerate}

Drawing a card from $\Pile_{\text{draw}}$ always ends the player's turn
(except when answering a question). Playing a card may trigger special effects
before the turn advances.

\subsection{Draw Pile Exhaustion}

If $\Pile_{\text{draw}}$ is empty when a player must draw, the discard pile
$\Pile_{\text{discard}}$ (excluding the current top card $\Top$) is shuffled
and becomes the new draw pile. $\Top$ remains as the sole card in
$\Pile_{\text{discard}}$.

% ══════════════════════════════════════════════════════════════════════════════
\section{Special Cards}
\label{sec:special}
% ══════════════════════════════════════════════════════════════════════════════

\subsection{Bombs (2, 3, Joker)}
\label{sec:bomb}

\begin{definition}[Bomb value]
  The bomb value $\func{bomb}(c)$ of a card is:
  \[
  \func{bomb}(c) = \begin{cases}
    2 & \text{if } c = \Card{s}{2} \\
    3 & \text{if } c = \Card{s}{3} \\
    5 & \text{if } c \text{ is a joker} \\
    0 & \text{otherwise}
  \end{cases}
  \]
\end{definition}

\begin{rule}[Initiating a bomb]
  When a player plays a bomb card (2, 3, or joker legally matching $\Top$),
  a bomb state $\Bomb$ is created targeting the next player in the current
  direction. The bomb count is initialized to $\func{bomb}(c)$.
\end{rule}

\begin{rule}[Bomb response]
  The targeted player must choose exactly one of:
  \begin{enumerate}[nosep, label=(\alph*)]
    \item \textbf{Pick.} Draw $\Bomb.\text{count}$ cards from the draw pile.
          The bomb is cleared and the player's turn ends.
    \item \textbf{Counter with bomb.} Play a bomb card (2, 3, or joker) from
          their hand. The bomb count increases by $\func{bomb}(c)$, and the
          bomb passes to the \emph{next} player in the current direction.
          \textbf{Bombs are cumulative and chainable.}
    \item \textbf{Counter with Ace.} Play any Ace from their hand. The bomb
          is cleared entirely (see~\S\ref{sec:ace}).
  \end{enumerate}
\end{rule}

\begin{example}
  Player A plays $\Card{\heartsuit}{2}$ (bomb 2) on $\Card{\heartsuit}{9}$.
  Player B is targeted.\\
  B counters with $\Card{\clubsuit}{3}$ (bomb 3) $\to$ cumulative bomb count $= 5$.\\
  Player C is now targeted. C picks 5 cards from the draw pile.
\end{example}

\subsection{Questions (8, Q)}
\label{sec:question}

\begin{rule}[Playing a question]
  When a player plays an 8 or Q legally matching $\Top$, they enter a
  \emph{question state}. Multiple 8's and Q's may be chained in a single
  turn, but \emph{only the first must be legal} against $\Top$; subsequent
  chained 8/Q cards must be legal against the card immediately preceding
  them.
\end{rule}

\begin{rule}[Answering]
  After the question chain, the \textbf{same player} must answer by playing
  a \emph{legal, non-special} card. If they cannot, they must draw one card
  from $\Pile_{\text{draw}}$. The question state is then cleared and the
  player's turn ends.
\end{rule}

\begin{note}
  The questioner answers their own question---the question does \emph{not}
  pass to the next player.
\end{note}

\subsection{Jump (J)}
\label{sec:jump}

\begin{rule}[Jump]
  When a player plays a J legally matching $\Top$, the \emph{next} player in
  the current direction is skipped. Multiple J's may be chained: $m$ J's skip
  $m$ players. The skipped player(s) may \textbf{reject} the jump by playing
  a J of their own; upon rejection, the rejecting player takes their turn
  immediately.
\end{rule}

\begin{example}
  Direction is clockwise. Player A plays JJ (2 J's).\\
  $\to$ Players B and C are skipped. Player D plays.\\
  \textbf{Rejection:} Player B plays a J of their own $\to$ B takes their
  turn immediately; C is not skipped.
\end{example}

\subsection{Kickback (K)}
\label{sec:kickback}

\begin{rule}[Kickback]
  When one or more K's are played:
  \begin{itemize}[nosep]
    \item An \textbf{odd} number of K's reverses the direction of play
          ($\Dir \leftarrow -\Dir$). Turn passes to the next player in the
          new direction.
    \item An \textbf{even} number of K's cancels out (direction unchanged).
          The playing player \textbf{continues} and may play another card.
  \end{itemize}
\end{rule}

\begin{example}
  Direction is clockwise. Player A plays K (odd $\to$ reverse).\\
  Direction is now counterclockwise. Next player is the one ``before'' A.\\
  \textbf{Even case:} A plays KK $\to$ direction unchanged, A plays again.
\end{example}

\subsection{Aces}
\label{sec:ace}

\subsubsection{Regular Aces ($\Card{s}{A}, s \neq \spadesuit$)}

\begin{rule}[Regular Ace --- defensive]
  When used as a counter against a bomb: the bomb is cleared. The Ace acts
  solely as a shield; no request power is granted.
\end{rule}

\begin{rule}[Regular Ace --- offensive]
  When played legally (matching $\Top$ by suit or rank): the playing player
  may request \textbf{either} a specific suit \textbf{or} a specific rank
  (not both). Subsequent players, in turn order, must either play a card
  matching the request or draw from $\Pile_{\text{draw}}$. The request
  persists until satisfied.
\end{rule}

\subsubsection{Ace of Spades ($\Card{\spadesuit}{A}$)}

The Ace of Spades wields \textbf{dual power}.

\begin{rule}[A$\spadesuit$ --- defensive]
  When used as a counter against a bomb: one power is consumed to stop the
  bomb. The remaining power allows the player to request \textbf{either} a
  suit \textbf{or} a rank (not both).
\end{rule}

\begin{rule}[A$\spadesuit$ --- offensive]
  When played legally (matching $\Top$ by suit or rank): \textbf{both} powers
  are available. The player may request \textbf{both} a suit \textbf{and} a
  rank simultaneously (effectively requesting a specific card, excluding
  jokers).
\end{rule}

\begin{rule}[Ace request resolution]
  When an Ace request $\CardReq$ is active:
  \begin{itemize}[nosep]
    \item If $\CardReq$ specifies both suit $s$ and rank $r$, the next player
          must play $\Card{s}{r}$ or draw.
    \item If $\CardReq$ specifies only suit $s$, the next player must play any
          $\Card{s}{\cdot}$ or draw.
    \item If $\CardReq$ specifies only rank $r$, the next player must play any
          $\Card{\cdot}{r}$ or draw.
  \end{itemize}
  If a player draws, $\CardReq$ passes to the next player. This continues
  until someone plays the requested card, at which point $\CardReq$ is
  cleared. If the draw pile runs out, the request ends.
\end{rule}

% ══════════════════════════════════════════════════════════════════════════════
\section{The ``Card'' Declaration}
\label{sec:declare}
% ══════════════════════════════════════════════════════════════════════════════

This rule is the Kenyan Poker analogue of UNO's ``UNO!'' call.

\begin{definition}[On the verge of winning]
  A player is \textbf{on the verge of winning} iff, after their play, their
  remaining hand has at most 2 cards and contains at least one non-special
  card. Formally, player $p$ is on the verge iff:
  \[
  0 < |\Hand_p| \leq 2 \;\land\; \exists c \in \Hand_p : c \notin \Special
  \]
\end{definition}

\begin{rule}[Declaration obligation]
  Immediately after playing a card that leaves them on the verge, a player
  \textbf{must declare ``Card!''} Failure to declare results in:
  \begin{itemize}[nosep]
    \item The player is flagged as having \emph{failed to declare}
          ($p \in \Failed$).
    \item Should they subsequently play what would have been a winning
          (non-special) last card, they are \textbf{penalized}: they do
          \emph{not} win, and must instead draw one card from
          $\Pile_{\text{draw}}$.
    \item The failure flag is cleared after the penalty is applied.
  \end{itemize}
\end{rule}

\begin{remark}[Strategic implications]
  When a player is flagged as on the verge ($p \in \OnVerge$), all players
  are aware of their imminent victory. This enables counterplay:
  \begin{itemize}[nosep]
    \item Use K (kickback) to reverse direction so a bomb-holding player can
          bomb the verge player before they win.
    \item Use J (jump) to skip the verge player.
    \item Chain bombs onto the verge player to force them to pick cards,
          pushing them off the verge.
  \end{itemize}
\end{remark}

% ══════════════════════════════════════════════════════════════════════════════
\section{Win Conditions and Cardless State}
\label{sec:win}
% ══════════════════════════════════════════════════════════════════════════════

\begin{definition}[Winning]
  A player $p$ wins iff:
  \begin{enumerate}[nosep, label=(\roman*)]
    \item They play a card $c$ (legally matching $\Top$),
    \item $c$ is \textbf{non-special} ($c \notin \Special$),
    \item $\Hand_p = \varnothing$ after the play,
    \item $p \notin \Failed$ (they did not fail a ``card'' declaration).
  \end{enumerate}
  The game ends when $n-1$ players have won (the last remaining player loses).
\end{definition}

\begin{definition}[Cardless state]
  \label{sec:cardless}
  If a player plays a \textbf{special} card as their last card
  ($\Hand_p = \varnothing$ and $c \in \Special$), they enter the
  \textbf{cardless} state ($p \in \Cardless$):
  \begin{itemize}[nosep]
    \item No player can win during a round where any player is cardless.
    \item On the cardless player's next turn, they \textbf{must} draw one
          card from $\Pile_{\text{draw}}$. The cardless flag is cleared, and
          their turn ends.
    \item After all cardless players have drawn, normal play and winning
          resume.
  \end{itemize}
\end{definition}

% ══════════════════════════════════════════════════════════════════════════════
\section{Formal Semantics}
\label{sec:formal}
% ══════════════════════════════════════════════════════════════════════════════

This section provides a mathematical specification of the game as a transition
system. It serves as the single source of truth; the TypeScript implementation
is a direct encoding of these definitions.

\subsection{Game State}

\begin{definition}[Game state]
  The game state $\State$ is a tuple:
  \[
  \State = (\Players, \Pile_{\text{draw}}, \Pile_{\text{discard}}, \Top, \Dir, i_{\text{curr}},
            \Bomb, \Quest, \CardReq, \OnVerge, \Failed, \Winner, \Cardless,
            \text{gameOver})
  \]
  where:
  \begin{itemize}[nosep]
    \item $\Players = (p_0, p_1, \dots, p_{n-1})$ is the ordered list of
          players, each $p_k = (id, \Hand_k)$.
    \item $\Pile_{\text{draw}}$ is the draw pile (a sequence of cards).
    \item $\Pile_{\text{discard}}$ is the discard pile.
    \item $\Top$ is the current top card.
    \item $\Dir \in \{+1, -1\}$ is the direction of play.
    \item $i_{\text{curr}} \in \{0, \dots, n-1\}$ is the current player index.
    \item $\Bomb \in \{\bot\} \cup \{(c, o)\}$ where $c \in \mathbb{N}^+$ is
          the cumulative bomb count and $o$ is the originating player.
    \item $\Quest \in \{\bot\} \cup \{(m, q)\}$ where $m$ is the chain count
          and $q$ is the questioning player.
    \item $\CardReq \in \{\bot\} \cup \{(s, r)\}$ where $s$ is an optional
          suit and $r$ is an optional rank.
    \item $\OnVerge, \Failed \subseteq \Players$ are sets of player IDs.
    \item $\Winner \subseteq \Players$ and $\Cardless \subseteq \Players$.
    \item $\text{gameOver} \in \{\top, \bot\}$.
  \end{itemize}
\end{definition}

\subsection{Transition Function}

Let $\func{step} : \State \times \Action \to \State$ be the state transition
function. We define the most important transitions below.

\subsubsection{Normal Play}

When player $p = \Players[i_{\text{curr}}]$ plays card $c$ at hand index
$j$ on top card $\Top$:

\begin{enumerate}[nosep]
  \item \textbf{Precondition:} $\func{legal}(c, \Top) = \top$.
  \item \textbf{Effect:}
        $\Hand_p \leftarrow \Hand_p \setminus \{c\}$;
        $\Pile_{\text{discard}} \leftarrow \Pile_{\text{discard}} \cdot c$;
        $\Top \leftarrow c$.
  \item \textbf{Win check:} If $\Hand_p = \varnothing \land c \notin \Special
        \land p \notin \Failed$, then $\Winner \leftarrow \Winner \cup \{p\}$.
  \item \textbf{Cardless check:} If $\Hand_p = \varnothing \land c \in \Special$,
        then $\Cardless \leftarrow \Cardless \cup \{p\}$.
  \item \textbf{Verge check:} If $0 < |\Hand_p| \leq 2 \land \Hand_p \cap \Normal
        \neq \varnothing$, then $\OnVerge \leftarrow \OnVerge \cup \{p\}$.
        If the accompanying action had $\text{declareCard} = \bot$, then
        $\Failed \leftarrow \Failed \cup \{p\}$.
  \item \textbf{Bomb effect:} If $c$ is a bomb card, $\Bomb \leftarrow
        (\func{bomb}(c), p)$.
  \item \textbf{Question effect:} If $c$ is 8 or Q, $\Quest \leftarrow (m, p)$.
  \item \textbf{Jump effect:} If $c$ is J, advance $i_{\text{curr}}$ past
        $m$ players.
  \item \textbf{Kickback effect:} If $c$ is K, apply direction logic.
  \item \textbf{Ace effect:} If $c$ is an Ace, set $\CardReq$ accordingly.
  \item \textbf{Turn advance:} $i_{\text{curr}} \leftarrow
        (i_{\text{curr}} + \Dir) \bmod n$.
\end{enumerate}

\subsubsection{Bomb Counter}

When a player counters with bomb card $c$:
\[
\Bomb \leftarrow (\Bomb.\text{count} + \func{bomb}(c), \Bomb.\text{origin})
\]
The bomb passes to the next player.

\subsubsection{Bomb Clear with Ace}

When a player plays an Ace as a bomb counter:
\[
\Bomb \leftarrow \bot
\]
If $c = \Card{\spadesuit}{A}$, the player may set $\CardReq$ with one power
(either suit or rank).

\subsubsection{Draw from Pile}

When a player draws:
\[
\Hand_p \leftarrow \Hand_p \cdot \func{draw}(1),\quad
\Pile_{\text{draw}} \leftarrow \Pile_{\text{draw}} \setminus \{\text{top card}\}
\]
If $\Pile_{\text{draw}} = \varnothing$, recycle $\Pile_{\text{discard}}$
(excluding $\Top$).

\subsection{Turn Priority}

The transition function dispatches to the appropriate sub-transition based
on the current state, with the following strict priority:

\[
\text{Active}(p) = \begin{cases}
  \text{cardless}   & \text{if } p \in \Cardless \\
  \text{bomb}       & \text{if } \Bomb \neq \bot \\
  \text{question}   & \text{if } \Quest \neq \bot \land \Quest.q = p \\
  \text{ace\_req}   & \text{if } \CardReq \neq \bot \\
  \text{normal}     & \text{otherwise}
\end{cases}
\]

% ══════════════════════════════════════════════════════════════════════════════
\section{Invariants and Properties}
\label{sec:invariants}
% ══════════════════════════════════════════════════════════════════════════════

\begin{invariant}[Card conservation]
  At all times, the total number of cards in the game is constant:
  \[
  |\Pile_{\text{draw}}| + |\Pile_{\text{discard}}| + \sum_{p \in \Players} |\Hand_p| = k \cdot 54
  \]
  where $k$ is the number of decks.
\end{invariant}

\begin{invariant}[Winner validity]
  $\forall p \in \Winner : p \notin \Failed$, i.e., no player who failed to
  declare ``card'' can appear in the winner set.
\end{invariant}

\begin{invariant}[Game-over condition]
  $\text{gameOver} = \top \iff |\Winner| = |\Players| - 1$.
\end{invariant}

\begin{invariant}[Cardless blocking]
  If $\Cardless \neq \varnothing$, then no new winners can be added to
  $\Winner$ during the current round---only after all cardless players have
  drawn.
\end{invariant}

\begin{proposition}[Termination]
  Assuming finite decks and the draw-pile recycling mechanism, every game of
  Kenyan Poker terminates in a finite number of turns.
\end{proposition}

\begin{proof}[Sketch]
  Each turn either (a) removes a card from a player's hand, (b) transfers
  cards between draw pile and hands, or (c) enforces a player to draw. Cards
  only move between hands and piles---none are created. The bomb and question
  mechanics are bounded: bomb chains must terminate (a player eventually picks
  or counters with an Ace), and question chains end when the questioner
  answers or draws. Since winning removes a player from the active set and
  $n-1$ wins end the game, the game terminates.
\end{proof}

% ══════════════════════════════════════════════════════════════════════════════
\section{Glossary of Notation}
\label{sec:notation}
% ══════════════════════════════════════════════════════════════════════════════

\begin{longtable}{c|l}
\toprule
\textbf{Symbol} & \textbf{Meaning} \\
\midrule
\endfirsthead
\midrule
\textbf{Symbol} & \textbf{Meaning} \\
\midrule
\endhead
$\Card{s}{r}$     & Regular card of suit $s$, rank $r$ \\
$\Joker{col}$     & Joker of color $col \in \{\text{red}, \text{black}\}$ \\
$\Special$        & Set of all special cards \\
$\Normal$         & Set of all non-special (normal) cards \\
$\Deck$           & A standard 54-card deck \\
$\Hand_p$         & Hand of player $p$ \\
$\Pile_{\text{draw}}$    & Draw pile (face-down) \\
$\Pile_{\text{discard}}$ & Discard pile (face-up) \\
$\Top$            & Top card of the discard pile \\
$\Dir$            & Direction of play ($+1$ or $-1$) \\
$\Players$        & Ordered list of players \\
$i_{\text{curr}}$ & Index of the current player \\
$\Bomb$           & Bomb state: $(count, origin)$ or $\bot$ \\
$\Quest$          & Question state: $(chain\_count, player)$ or $\bot$ \\
$\CardReq$        & Ace request: $(suit?, rank?)$ or $\bot$ \\
$\OnVerge$        & Set of players who must declare ``card'' \\
$\Failed$         & Set of players who failed to declare \\
$\Winner$         & Set of players who have won \\
$\Cardless$       & Set of players who are cardless \\
$\func{bomb}(c)$  & Bomb value of card $c$ \\
$\func{color}(c)$ & Color of card $c$ \\
$\func{legal}(c, \Top)$ & Whether $c$ is legal on top of $\Top$ \\
\bottomrule
\end{longtable}

\end{document}
